% !TeX spellcheck = ru_RU
% !TEX root = vkr.tex

\section{Обзор}
\label{sec:related}

\subsection{Кодирование Хаффмана}

Алгоритм Хаффмана строит оптимальный префиксный код для известного
распределения частот символов~\cite{huffman1952}. Сначала для каждого символа
создаётся лист дерева с весом, равным его частоте. Затем два узла с
минимальными весами объединяются во внутренний узел с весом, равным сумме весов
объединённых узлов; операция повторяется, пока не останется один корень. Путь
от корня к листу задаёт двоичный код символа: переход влево соответствует биту
0, переход вправо~--- биту 1.

Префиксность означает, что ни один код символа не является началом другого
кода. Поэтому декодирование выполняется последовательным обходом дерева по
битам входного потока: достижение листа однозначно определяет восстановленный
символ. В статическом варианте кодирования распределение частот определяется
до начала записи сжатых данных; декодеру необходимо получить сведения,
позволяющие восстановить то же дерево.

\subsection{Существующие решения}

В алгоритме \Deflate{} кодирование Хаффмана применяется после этапа
LZ77, который заменяет повторяющиеся фрагменты ссылками на уже обработанные
данные~\cite{rfc1951}. Такой подход используется библиотекой zlib и
утилитой gzip~\cite{zlib,gzip}. Его преимуществом является более эффективное
сжатие типичных файлов, поскольку учитываются не только частоты отдельных
значений, но и повторяющиеся последовательности. Недостаток для целей данной
работы состоит в том, что результат измерений характеризует комбинацию
алгоритмов, а не отдельное кодирование Хаффмана.

Статический архиватор, реализованный в работе, занимает другую позицию: он
хранит таблицу частот и применяет только побайтовое кодирование Хаффмана. Он
уступает универсальным промышленным архиваторам как средство достижения
максимального коэффициента сжатия, но позволяет непосредственно изучать
зависимость результата от распределения байтов и сохраняет простой формат
архива.

Кроме промышленных архиваторов, существуют учебные и справочные реализации
кодирования Хаффмана. Например, Rosetta Code содержит набор реализаций
алгоритма на разных языках программирования~\cite{rosettahuffman}, а проект
Nayuki Reference Huffman coding предоставляет компактную справочную реализацию
кодирования и декодирования~\cite{nayukihuffman}. Такие материалы полезны как
иллюстрация алгоритма, но обычно не ставят акцент на переносимом бинарном
формате архива, обработке граничных случаев и воспроизводимом эксперименте. В
данной работе эти аспекты выделены как самостоятельные инженерные требования.

\subsection{Выводы}

Для реализации выбран статический вариант алгоритма Хаффмана. Он соответствует
цели работы, поскольку обеспечивает декодирование произвольных бинарных
данных, допускает явное описание формата архива и позволяет измерить свойства
именно энтропийного кодирования без влияния словарных методов.
